\documentclass[10pt,a4paper]{amsart}
\usepackage[margin=19mm]{geometry}
\usepackage{amsmath,amssymb,mathtools}
\usepackage[scr=rsfs,cal=euler]{mathalfa}
\usepackage{xfrac}
\usepackage[unicode,hidelinks,bookmarks=false]{hyperref}
\newcommand{\ie}{i.e.\ }
\newcommand{\cf}{cf.\ }
\newcommand{\resp}{respectively\ }
\newcommand{\Subsection}{\subsection}
\providecommand{\nobreakdash}{\nobreak\hbox{-}}
\setlength{\emergencystretch}{2em}
\multlinegap0pt
\usepackage{amssymb}
\usepackage{dsfont}
\usepackage{bm}
\usepackage{mathtools}
\usepackage[shortlabels]{enumitem}
\newtheorem{lemma}{Lemma}
\newcommand{\Cspace}{\mathbb{C}}
\newcommand{\dy}{\dot{\vy}}
\newcommand{\ex}{\re} % Euler's number
\newcommand{\mo}{\vPh_T}
\newcommand{\pd}[2]{\frac{\partial #1}{\partial #2}}
\newcommand{\re}{\mathrm e}
\newcommand{\rp}{\mathrm p}
\newcommand{\vA}{\mathbf A}
\newcommand{\vI}{\mathbf I}
\newcommand{\vJ}{\mathbf J}
\newcommand{\vP}{\mathbf P}
\newcommand{\vPh}{\mathbf \Phi}
\newcommand{\vPs}{\mathbf \Psi}
\newcommand{\vQ}{\mathbf Q}
\newcommand{\vV}{\mathbf V}
\newcommand{\vf}{\mathbf f}
\newcommand{\vx}{\mathbf x}
\newcommand{\vy}{\mathbf y}
\newcommand{\vzero}{\mathbf{0}}
\title{Koopman-based stability analysis of differential-algebraic equations with applications to frictional multibody systems}
\author{Alexander Sch\"utz, Fabia Bayer, Remco I. Leine}
\date{Source: arXiv:2607.17339v1 (CC BY 4.0)}
\begin{document}
\maketitle

\noindent\emph{Source and licence.} Koopman-based stability analysis of differential-algebraic equations with applications to frictional multibody systems. Version arXiv:2607.17339v1 (CC BY 4.0), distributed by arXiv under the Creative Commons Attribution 4.0 International licence, http://creativecommons.org/licenses/by/4.0/, as recorded in the arXiv licence metadata for this exact version and shown on the abstract page https://arxiv.org/abs/2607.17339v1. Full licence notice: \texttt{licenses/arxiv-2607.17339-CC-BY-4.0.txt}.\par\smallskip

\noindent\textit{Editorial scope.} Counted unit: the article's lemma eq:Floquetproof:Floquettheorem (source0.tex lines 1489--1494). The article's complete original proof of it is delivered with it (source0.tex lines 1496--1545, one \texttt{proof} environment of 50 lines). No mathematics is written, repaired or re-derived: the statement and the proof are the article's own lines, in the article's own notation, and no retained line is substituted or re-flowed.  Retained uncounted as the source-local context the proof needs: lines 726--728; lines 746--750.  Nothing was omitted from any retained span, so the package carries no omission record.  No retained line contains a citation, so the wrapper carries no bibliography and no bibliography entry.  No retained line contains a figure environment, an image-inclusion command or drawing code, so no image is delivered and the package declares no asset dependency.  Editorial formatting only, in addition: an AMS \texttt{amsart} preamble that restates the article's own declarations and macros used by the retained lines, namely the environments \texttt{lemma} on separate counters, together with the standard packages those lines need.  The retained environments print in this document as Lemma 1 in order of appearance, whereas the article prints the same items with the numbering of its own full document.  The complete native source of the article as distributed in the arXiv e-print for this exact version is bundled unchanged as \texttt{sources/205562/source0.tex}.
\begin{equation}\label{eq:KHDAE:LTPDAE}
    \vA\dy = \vJ(t)\vy\textnormal{,}\quad \vJ(t+T) = \vJ(t):= \left.\pd{\vf}{\vx}\right|_{\vx(t)=\vx^\rp(t)}\textnormal{.}
\end{equation}

\begin{align}\label{eq:dae:fundamat}
    \vPh(t, t_0) := \vV(t) \begin{pmatrix}
        \vPs(t) \vPs(t_0)^{-1} & \vzero \\ \vzero & \vzero
    \end{pmatrix} \vV(t_0)^{-1} \;.
\end{align}

\begin{lemma}[Floquet's theorem for DAEs]\label{eq:Floquetproof:Floquettheorem}
    Consider the linear time-periodic DAE~\eqref{eq:KHDAE:LTPDAE} with constant solution space dimension $r$ and its fundamental solution matrix given in Equation~\eqref{eq:dae:fundamat}. There exists a $T$-periodic matrix $\hat{\vP}(t) \in \Cspace^{r \times r}$ and a constant matrix $\vQ \in \Cspace^{r \times r}$ such that
    \begin{align}
        \vPs(t) \vPs(0)^{-1} = \hat{\vP}(t) \ex^{\vQ t} \;.
    \end{align}
\end{lemma}

\begin{proof}
    Due to the periodicity of the DAE~\eqref{eq:KHDAE:LTPDAE}, the fundamental solution matrix has a $T$-shift property in both arguments simultaneously, i.e., 
    \begin{align}\label{eq:Floquetproof:periodicity}
        \vPh(t, 0) = \vPh(t+T, T), \quad t \geq 0 \;.
    \end{align}
    As discussed above, the group property of the fundamental solution matrix holds in forward time, i.e., 
    \begin{align}\label{eq:Floquetproof:group1}
        \vPh(t + T, 0) = \vPh(t + T, T) \vPh(T, 0), \quad t \geq 0 \;.
    \end{align}
    Combining \eqref{eq:Floquetproof:periodicity} and \eqref{eq:Floquetproof:group1} yields
    \begin{align}\label{eq:FloquetProof:group}
        \vPh(t + T, 0) =  \vPh(t, 0) \mo \;.
    \end{align} 
    Since the matrix $\vPs(t) \vPs(0)^{-1}$ is invertible for all $t \geq 0$, there exists a matrix $\vQ \in \Cspace^{r \times r}$ such that 
    \begin{align}\label{eq:FloquetProof:matrixlog}
        \ex^{\vQ T} = \vPs(T) \vPs(0)^{-1}\;.
    \end{align}
The matrix $\hat{\vP}(t) := \vPs(t) \vPs(0)^{-1} \ex^{-\vQ t}$ is then defined appropriately such that
    \begin{align}\label{eq:FloquetProof:PhiP}
        \vPh(t,0) = \vV(t) \begin{pmatrix}
            \hat{\vP}(t) \ex^{\vQ t} & \vzero\\ \vzero & \vzero
        \end{pmatrix} \vV(0)^{-1}
    \end{align}
    holds. Clearly, $\hat{\vP}(T)$ is the identity matrix. It remains to show that $\hat{\vP}(t)$ is indeed $T$-periodic. To this end, Equation~\eqref{eq:FloquetProof:PhiP} is evaluated at time $t+T$ as
    \begin{align}\label{eq:FloquetProof:PhitplusT}
        \vPh(t + T,0) = \vV(t) \begin{pmatrix}
            \hat{\vP}(t + T) \ex^{\vQ (t+T)}& \vzero\\ \vzero & \vzero
        \end{pmatrix} \vV(0)^{-1}\;.
    \end{align}
    On the other hand, substituting the definitions of $\hat{\vP}(t)$ and $\vQ$ into Equation~\eqref{eq:FloquetProof:group} yields
    \begin{align}\label{eq:FloquetProof:PhitplusT2}
        &\vPh(t + T, 0) =  \vPh(t, 0) \mo \nonumber\\
        =& \vV(t) \begin{pmatrix}
            \hat{\vP}(t) \ex^{\vQ t} & \vzero\\ \vzero & \vzero
        \end{pmatrix} \begin{pmatrix}
            \hat{\vP}(T)\ex^{\vQ T} & \vzero\\ \vzero & \vzero
        \end{pmatrix} \vV(0)^{-1}
    \end{align}
    by making use of the property $\vV(0)^{-1}\vV(T) =\vI$ due to the $T$-periodicity of the matrix $\vV(t)$.
    
    By equating the expressions \eqref{eq:FloquetProof:PhitplusT} and \eqref{eq:FloquetProof:PhitplusT2}, acknowledging the property $\hat{\vP}(T)=\vI$, pre-multiplying with $\vV(t)^{-1}$, and post-multiplying with $\vV(0)$, we obtain
\begin{align}
        \begin{pmatrix}
            \hat{\vP}(t + T) \ex^{\vQ (t+T)} & \vzero\\ \vzero & \vzero
        \end{pmatrix} = \begin{pmatrix}
            \hat{\vP}(t) \ex^{\vQ (t+T)} & \vzero\\ \vzero & \vzero
        \end{pmatrix} \;.
    \end{align}
    With $\ex^{\vQ (t+T)}$ invertible, this directly implies the desired periodicity property of the matrix $\hat{\vP}(t)$.
\end{proof}

\end{document}
